\section{Upper bounds at the prescribed radius}
\label{sec:upper}

For a set \(T\) of labels, define its overlap and the directed deficit
on its complement by
\[
 O_T(p,q)=\sum_{i\in T}\min(p_i,q_i),\qquad
 L_H(p,q)=\sum_{i\in H}(q_i-p_i)_+,\qquad H=T^\complement .
\]
Writing \(w=q(T)\), normalization gives
\begin{equation}
 \operatorname{TV}(p,q)=w+L_H(p,q)-O_T(p,q).
 \label{eq:upper-head-tail}
\end{equation}
In particular, at \(\Delta=e-w\), the strict alternative is exactly
\(L_H(p,q)>\Delta+O_T(p,q)\).

\subsection{A capped overlap test}

For a positive cap \(b\), replace each selected label \(i\) by
\(m_i=\max(1,\lceil q_i/b\rceil)\) equiprobable copies. Use independent
copy choices under both input laws. On the refined labels,
\(p_i\) and \(q_i\) are divided by the same \(m_i\). The total mass and
overlap on any union of full copy classes are preserved, and every refined
reference probability in the selected set is at most \(b\). Its refined
moment satisfies
\begin{equation}
 \frac{1}{\sqrt{2}}\sum_{i\in T}\min(q_i\sqrt b,q_i^{3/2})
 \ \le\ \sum_{i\in T}\frac{q_i^{3/2}}{\sqrt{m_i}}
 \ \le\ \sum_{i\in T}\min(q_i\sqrt b,q_i^{3/2}).
 \label{eq:upper-clone-moment}
\end{equation}
For \(q_i\le b\) the atom is unchanged. For \(q_i>b\) use
\(q_i/b\le m_i<2q_i/b\). Null-zero labels have one copy and remain
available to the alternative.

We now work on a refined alphabet and denote its probabilities by \(p_i,q_i\).
Suppose \(q_i\le b\) for \(i\in T\), \(0<b\le x\), and
\[
   M=\sum_{i\in T}q_i^{3/2}\ge 8x\sqrt b.
\]
We test against all full laws with \(O_T(p,q)\le x\), with no condition
on the number of tail observations or on \(p(T)\). Put
\(\ell=\lceil8192\sqrt{x}/M\rceil\) and \(r=M^2/x\).
On a first batch of \(\ell\) full observations use the score
\(W(i)=\sqrt{q_i}\) on \(T\) and \(W(i)=0\) off \(T\). Reject if its empirical
mean is below \(7M/8\). Under any law,
\(\operatorname{Var}(W)\le\sqrt b\,\mathbb E W\).
Thus this gate has null error at most \(64\sqrt b/(\ell M)<1/100\),
and it misses an alternative of mean at most \(3M/4\) with probability
at most the same bound.

If instead the mean exceeds \(3M/4\), let
\(F=\{i\in T:p_i\ge q_i\}\). The score contribution from \(T\setminus F\)
is at most \(\sqrt b\,O_T\le M/8\), and \(q(F)\le O_T\le x\).
Cauchy--Schwarz on \(F\) shows that
\[
  \sum_{i\in T}p_i^2\ge M^2/(4x)=r/4.
\]
The null squared mass is at most \(\sqrt b\,M\le r/8\).

In two new independent batches of size \(\ell\), count matching pairs
with one observation from each batch and with label in \(T\). Reject if
the cross-pair proportion exceeds \(3r/16\). For any law \(p\), put
\(B=\sum_{i\in T}p_i^2\) and \(C=\sum_{i\in T}p_i^3\). The cross count
has mean \(\ell^2 B\). Two indicators for disjoint observation positions
have covariance zero. Indicators that share exactly one position
have joint expectation \(C\), and the identical pair has expectation
\(B\). Counting ordered indicator pairs gives the variance bound
\[
  \operatorname{Var}(\text{cross count})
  \le \ell^2 B+2\ell^3 C
  \le \ell^2 B+2\ell^3 B^{3/2}.
\]
In the last step, every \(p_i\le\sqrt{B}\) on \(T\).
By Chebyshev, the cross-pair gate has null error at most
\(32/(\ell^2 r)+64/(\ell\sqrt r)\), and miss probability
at most \(16/(\ell^2 B)+32/(\ell\sqrt B)\) when \(B\ge r/4\).
Both are below \(1/100\) for \(\ell\sqrt r\ge8192\).
The combined test on \(3\ell\) draws therefore has null error at most
\(1/50\) and miss probability at most \(1/100\).
This is the standard collision moment argument on a selected
subprobability, with a separate weighted score to enforce power
for arbitrary nonuniform alternatives \cite{paninski2008}.

\subsection{Close to the maximum distance}

Let \(\delta=1-e>0\) and
\[
 J=\sum_i q_i\min(1,\sqrt{q_i/\delta}).
\]
When \(J\ge16\delta\), split all labels at cap \(b=\delta\). By
\eqref{eq:upper-clone-moment} the refined moment is at least
\(\sqrt{\delta}\,J/\sqrt 2\ge8\delta\sqrt\delta\).
A strict alternative has full overlap below \(\delta\).
The capped overlap test with \(x=\delta\) uses
\(3\lceil16384/J\rceil\) observations or fewer, including the randomized
refinement. This proves the upper bound \(O(1/J)\) in this branch.

\subsection{The regular regime}

Use the shortest descending prefix \(H\) with \(q(H)\ge\delta\) and write
\(A=q(H)=1-w\). Choose \(i_*\) to be its largest-probability label and put
\(O=H\setminus\{i_*\}\). The head term in original probability units is
\[
 P(n)=\min_{\substack{R\ \mathrm{prefix}\\\mathrm{of}\ O}}
  \left\{\frac{\big(\sum_{i\in R}q_i^{2/3}\big)^{3/4}}{\sqrt n}
              +q(O\setminus R)\right\},\qquad
 D(n)=\sqrt{wA/n}+\frac1n+P(n).
\]
For the normalized conditional head \(h=q_H/A\), this is
\(P(n)=A P_h(nA)\) by direct scaling. No conditional observations are
required for the test below.

First test the tail mass using \(n\) full observations. For the tail count,
use the threshold \(512L\) about \(nw\), where \(L=\sqrt{nwA}+1\).
Its null variance is \(nwA\), giving null error less than \(1/100\).
At alternative mass \(s\), the variance is at most \(n(wA+|s-w|)\),
since \(s(1-s)-w(1-w)=(s-w)(1-s-w)\).
If \(n|s-w|>1024L\), accepting requires a deviation from the alternative
mean of at least \(n|s-w|/2\). Writing \(d=|s-w|\), the resulting miss bound is
\(4(nwA+nd)/(nd)^2\le4/1024^2+4/1024<1/100\).
In particular the same test detects
\(p(S)>4096(q(S)+1/n)\) for an arbitrary prescribed set \(S\), on
replacing \(T\) by \(S\). To see the implication, put \(u=nq(S)\) and
use \(\sqrt{u}+1\le u+2\).

For the remaining head shape, choose a minimizing prefix \(R\) and set
\[
 W_R=\sum_{i\in R}q_i^{2/3},\qquad B_R=W_R^{3/4}/\sqrt n,
 \qquad v=\sqrt{W_R}/n.
\]
If \(W_R=0\), take the retained set \(K\) empty. Otherwise let
\(K=\{i\in R:n^2v q_i^{2/3}\ge1\}\). The extra mass omitted by this
trimming is at most \(W_R/(n\sqrt v)=B_R\).
Take weights \(u_i=q_i^{-2/3}\) on \(K\), zero elsewhere, and define
\[
 Z=\sum_i u_i (N_i-nq_i)(N'_i-nq_i),\qquad
 V=\sum_i u_i(p_i-q_i)^2,
\]
from independent \(n\)-sample histograms \(N,N'\) of the full law.
This is the weighted centered quadratic method used in identity
testing \cite{valiantvaliant2017,chhorcarpentier2022}, here in a fixed
two-batch form. Its expectation is \(n^2 V\). Writing
\(\Sigma=\operatorname{diag}(p)-pp^\top\) for the single-observation
covariance, the centered bilinear part has variance
\(n^2\sum_{i,j}u_i u_j\Sigma_{ij}^2\le2n^2\sum_i u_i^2p_i^2\).
Indeed the diagonal term is at most \(\sum_i u_i^2p_i^2\) and the
off-diagonal term is at most \((\sum_i u_i p_i^2)^2\), which is bounded
by the same quantity. The linear parts are uncorrelated with the
bilinear part and have combined variance at most
\(2n^3\sum_i u_i^2p_i(p_i-q_i)^2\). In particular
\[
 \operatorname{Var}(Z)
 \le 3n^2\sum_i u_i^2p_i^2+
       6n^3\sum_i u_i^2p_i(p_i-q_i)^2.
\]
The cap on \(K\) gives \(\max u_i\le n^2v\), while
\(\sum_i u_i^2q_i^2\le W_R=n^2v^2\) and \(\max u_iq_i\le nv\).
It follows that
\[
 \sum_i u_i^2p_i^2\le2W_R+2n^2vV,\qquad
 \max_i u_i p_i\le n v+n\sqrt{vV}.
\]
Rejecting for \(Z>32n^2v\) gives null error at most \(3/32^2\).
If \(\sum_{i\in K}|p_i-q_i|>10000B_R\), Cauchy--Schwarz and
\(B_R^2=W_Rv\) give \(V>10^8v\). The threshold is below half the
mean. The ratio of variance to squared mean is bounded by
\[
 6(v/V)^2+12(v/V)+6\sqrt{v/V}.
\]
Four times this bound is below \(1/100\), giving the stated miss
probability by Chebyshev.

Apply the mass gate on another \(n\) observations to \(O\setminus K\).
Its null mass is at most \(B_R+q(O\setminus R)=P(n)\).
Unless one of the two shape guarantees applies, we have
\[
 \sum_{i\in O}|p_i-q_i|
 \le 14097P(n)+4096/n .
\]
For \(W_R=0\) the retained gate is unnecessary and the same bound holds.
Using \(|p(T)-w|+\sum_{i\in O}|p_i-q_i|\) to bound \(L_H(p,q)\)
follows by eliminating the single coordinate \(a\) with normalization.
Consequently any strict alternative outside the head rejection guarantees
has \(L_H(p,q)\le G\), where
\[
 G=1024\left(\sqrt{wA/n}+1/n\right)+14097P(n)+4096/n
 \le 14097D(n).
\]

Write \(Y(n)=\Delta+X(n)\), and suppose \(4000000D(n)\le Y(n)\).
Then \(256G\le Y(n)\). If \(X(n)\le\Delta\), this implies \(G\le\Delta\)
and every strict alternative is covered by a head guarantee.
Otherwise put \(x=X(n)/256\). Then \(G\le\Delta+x\), so by
\eqref{eq:upper-head-tail} all remaining strict alternatives have
\(O_T(p,q)<x\). Also \(nx\ge1\), since \(D(n)\ge1/n\) and
\(Y(n)<2X(n)\).

On the tail split at cap \(b=1/(n^2x)\), keeping the other labels
unchanged. The defining root gives
\[
 \sum_{i\in T}\min(q_i/x,nq_i^{3/2}/\sqrt x)\ge16
\]
by comparing each term at \(x=X(n)/256\) to its value at \(X(n)\).
Equation \eqref{eq:upper-clone-moment} then gives a refined moment
\(M\ge8\sqrt x/n=8x\sqrt b\).
Since \(b\le x\), the capped overlap test applies and needs at most
\(3\lceil8192\sqrt x/M\rceil\le3072n\) draws.
The group gate, head gates, and optional tail test use disjoint
batches of full observations. Combining their rejections gives both
errors below \(1/3\) with at most \(3076n\) observations. All inequalities
use the original \(e\), including equality at the population head gate.
